\documentclass[11pt,a4paper]{article}
\usepackage[utf8]{inputenc}
\usepackage[T1]{fontenc}
\usepackage{amsmath,amssymb,amsthm}
\usepackage{geometry}
\usepackage{enumitem}
\usepackage{hyperref}

\geometry{margin=1in}

\theoremstyle{plain}
\newtheorem{theorem}{Theorem}

\title{\textbf{The Collected Mathematical Papers of Arthur Cayley}\\[6pt]
\Large Volume V (1861--1863)\\[12pt]
\large Pages 76--100}
\author{Arthur Cayley}
\date{}

\usepackage{graphicx}
\emergencystretch=3em
\tolerance=600
\begin{document}
\maketitle
\thispagestyle{empty}
\newpage
\setcounter{page}{76}


% =====================================================================
% Continuation of Paper 315 (begun on p. 73), pp. 76 footer of paper
% =====================================================================

\section*{315. On the Cubic Centres of a Line with Respect to Three Lines and a Line \textit{(continued)}}

\noindent[Continuation from p.~75.]

\medskip

\noindent which is identically true in virtue of $a + b + c = 0$ (in fact, multiplying out, this gives
\begin{align*}
&12a^2b^2c^2 + 4(b^2c^4 + c^2a^4 + a^2b^4) + abc(a^6 + b^6 + c^6) \\
&\quad - 8a^2b^2c^2 - 4(b^2c^4 + c^2a^4 + a^2b^4) - 2abc(a^3 + b^3 + c^3) - a^3b^3c^3 = 0;
\end{align*}
that is
\[
3a^2b^2c^2 - abc(a^3 + b^3 + c^3) = 0, \quad \text{or} \quad abc(a^3 + b^3 + c^3 - 3abc) = 0,
\]
where the second factor divides by $a + b + c$), we find the above-mentioned equation.

We then have
\[
\frac{-x + y + z}{P} = \frac{x + y + z}{P} - \frac{2x}{P} = -3 + \frac{6a^2}{2a^2 + bc} = -\frac{3bc}{2a^2 + bc},
\]
that is
\[
\frac{-x + y + z}{P} = \frac{-3bc}{2a^2 + bc}, \quad \frac{x - y + z}{P} = \frac{-3ca}{2b^2 + ca}, \quad \frac{x + y - z}{P} = \frac{-3ab}{2c^2 + ab},
\]
and forming the product of these functions, and that of the foregoing values of $\frac{P}{x}$, $\frac{P}{y}$, $\frac{P}{z}$, we find, as before,
\[
-(-x + y + z)(x - y + z)(x + y - z) + xyz = 0
\]
for the equation of the locus of the single centre. The equation shows that the locus is a cubic curve which touches the lines $x = 0$, $y = 0$, $z = 0$ at the points where these lines are intersected by the lines $y - z = 0$, $z - x = 0$, $x - y = 0$ (that is, it touches the lines $x = 0$, $y = 0$, $z = 0$ harmonically in respect to the line $x + y + z = 0$), and besides meets the same lines $x = 0$, $y = 0$, $z = 0$ at the points in which they are respectively met by the line $x + y + z = 0$.

\medskip
\noindent 2, Stone Buildings, W.C., September 25, 1861.


\newpage
% =====================================================================
% Paper 316 — p. 77
% =====================================================================
\section*{316. Note on the Solution of an Equation of the Fifth Order.}
\addcontentsline{toc}{section}{316. Note on the Solution of an Equation of the Fifth Order}

[From the \textit{Philosophical Magazine}, vol.\ xxiii.\ (1862), pp.\ 195, 196.]

\medskip

This Note was in answer to Mr Jerrard's paper ``Remarks on Mr Cayley's Note,'' Phil.\ Mag.\ vol.\ xxi.\ pp.\ 348--350, referring to the foregoing paper 310.


\newpage
% =====================================================================
% Paper 317 — pp. 78–79
% =====================================================================
\section*{317. Note on the Transformation of a Certain Differential Equation.}
\addcontentsline{toc}{section}{317. Note on the Transformation of a Certain Differential Equation}

[From the \textit{Philosophical Magazine}, vol.\ xxiii.\ (1862), pp.\ 266, 267.]

\medskip

The differential equation
\[
(1 + x^2)\frac{d^2y}{dx^2} + x\frac{dy}{dx} - m^2 y = 0,
\]
if we put therein $i\theta = 2x^2 + 1$ ($i = \sqrt{-1}$ as usual), becomes
\[
(1 + \theta^2)\frac{d^2y}{d\theta^2} + \theta\frac{dy}{d\theta} - \tfrac{1}{4}m^2 y = 0.
\]

In fact an integral of the second equation is $(\sqrt{1+x^2} + x)^m$; this is
\[
= (\sqrt{(2x^2+1)^2 - 1} + 2x^2 + 1)^{\frac{1}{2}m};
\]
or putting $i\theta = 2x^2 + 1$, it is
\[
= \{\sqrt{-\theta^2 - 1} + i\theta\}^{\frac{1}{2}m};
\]
which is
\[
= \{i(\sqrt{\theta^2 + 1} + \theta)\}^{\frac{1}{2}m};
\]
so that an integral of the transformed equation in $\theta$ is
\[
(\sqrt{1+\theta^2} + \theta)^{\frac{1}{2}m};
\]
and writing in the second equation $\theta$ for $x$, and $\frac{1}{2}m$ for $m$, we see that the last-mentioned function, viz.\ $(\sqrt{\theta^2+1} + \theta)^{\frac{1}{2}m}$, is an integral of
\[
(1 + \theta^2)\frac{d^2y}{d\theta^2} + \theta\frac{dy}{d\theta} - \tfrac{1}{4}m^2 y = 0;
\]
\newpage
whence the transformed equation in $\theta$ must be this very equation, that is, it must be the first equation. I have for shortness used the particular integral $(\sqrt{1+x^2} + x)^m$; but the reasoning should have been applied, and it is in fact applicable, without alteration, to the general integral
\[
C(\sqrt{1+x^2} + x)^m + C'(\sqrt{1+x^2} - x)^m.
\]

There is of course no difficulty in a direct verification. Thus, starting from the first equation, or equation in $\theta$, the relation $i\theta = 2x^2 + 1$ gives
\[
\frac{dy}{d\theta} = \frac{i}{4x}\frac{dy}{dx}, \quad \frac{d^2y}{d\theta^2} = \frac{i}{4x}\frac{d}{dx}\!\left(\frac{i}{4x}\frac{dy}{dx}\right) = -\frac{1}{16x^2}\!\left(\frac{d^2y}{dx^2} - \frac{1}{x}\frac{dy}{dx}\right),
\]
\[
1 + \theta^2 = -4x^2(1+x^2);
\]
so that the equation becomes
\[
\tfrac{1}{4}(1+x^2)\!\left(\frac{d^2y}{dx^2} - \frac{1}{x}\frac{dy}{dx}\right) + \frac{1+2x^2}{4x}\frac{dy}{dx} - m^2 y = 0,
\]
or multiplying by 4,
\[
(1+x^2)\!\left(\frac{d^2y}{dx^2} - \frac{1}{x}\frac{dy}{dx}\right) + \frac{1+2x^2}{x}\frac{dy}{dx} - 4m^2 y = 0,
\]
that is
\[
(1+x^2)\frac{d^2y}{dx^2} + x\frac{dy}{dx} - 4m^2 y = 0,
\]
the second equation. But the first method shows the reason why the two forms are thus connected together.

\medskip
\noindent 2, Stone Buildings, W.C., February 19, 1862.


\newpage
% =====================================================================
% Paper 318 — pp. 80–84
% =====================================================================
\section*{318. On a Question in the Theory of Probabilities.}
\addcontentsline{toc}{section}{318. On a Question in the Theory of Probabilities}

[From the \textit{Philosophical Magazine}, vol.\ xxiii.\ (1862), pp.\ 361--365.]

\medskip

\textit{The question referred to is that discussed in the paper 121; the remarks on that paper in the Notes and References to volume II.\ are in a great measure to the same effect as the present and next papers, 318 and 319, the existence of which I had entirely overlooked. In the first part (dated 2, Stone Buildings, W.C., March 1862) of the present paper 318, after referring to the two modes of statement which may be called the Causation statement and the Concomitance statement, I reproduce nearly as in the Notes and References first my own solution as completed by Dedekind, next Boole's solution of the problem, involving his logical probabilities; and the paper is then continued as follows.}

\medskip

The foregoing paper was submitted to Prof.\ Boole, who, in a letter dated March 26, 1862, writes:

``The observations which have occurred to me after studying your paper are the following.

\medskip

``1st.\ I think that your solution is correct under conditions partly expressed and partly implied. The one to which you direct attention is the assumed independence of the causes denoted by $A$ and $B$. Now I am not sure that I can state precisely what the others are; but one at least appears to me to be the assumed independence of the events of which the probabilities according to your hypothesis are $\alpha\lambda$, $\beta\mu$. Assuming the independence of the causes as to happening, I do not think that you are entitled on that ground to assume their independence as to acting; because, to confine our observations to common experience, we often notice that states of things apparently independent as to their occurrence, may, when concurring, aid or hinder each other in such a manner that the one may be more or less likely to act `efficiently' in the presence of the other than in its absence. I use the language of your own hypothesis of efficient action.

``2ndly.\ When I say that I think your solution correct under certain conditions, I ought to add that it appears to me that such conditions ought to be stated as part of the original data, and that they ought to be of such a kind that they can be established by experience in the same way as the other data are. For instance, if experience, as embodied in a sufficiently long series of statistical records, establish that
\[
\text{Prob.\ } A = \alpha, \quad \text{Prob.\ } B = \beta,
\]
the very same experience may, by establishing also that
\[
\text{Prob.\ } AB = \alpha\beta,
\]
whence in conjunction with the former it follows that
\[
\text{Prob.\ } AB' = \alpha\beta', \quad \text{Prob.\ } A'B = \alpha'\beta, \quad \text{Prob.\ } A'B' = \alpha'\beta',
\]
enable us to pronounce that $A$ and $B$ are in the long run, as to happening or not happening, in the position of mutually independent events.

``3rdly.\ I think it may be shown to demonstration, from the nature of the result, that the solution you have obtained does not apply simply and generally to the problem under the single modification of the assumption that $A$ and $B$ are independent. The completed data under this assumption are
\begin{align*}
\text{Prob.\ } A &= \alpha, \quad \text{Prob.\ } B = \beta, \quad \text{Prob.\ } AB = \alpha\beta, \\
\text{Prob.\ } AE &= \alpha p, \quad \text{Prob.\ } BE = \beta q.
\end{align*}
You may deduce all these from your Table of Probabilities of `compound events' given in your paper. Now you may easily satisfy yourself that the sole necessary and sufficient conditions for the consistency of these data are the following:
\begin{align}
\alpha p' + \beta q &\geq \alpha\beta, \tag{1}\\
\alpha p + \beta q &\leq \alpha\beta, \tag{2}\\
\left\{\frac{\alpha}{\beta},\ \frac{\beta}{\alpha},\ \frac{p}{q},\ \frac{q}{p}\right\} &\leq 1. \tag{3}
\end{align}
\hfill (M).

But your solution requires the following conditions to be satisfied, viz.,
\[
q - \alpha p \geq 0, \quad p - \beta q \leq 0,
\]
together with the system (3). Now (1) and (2) are expressible in the form
\begin{align*}
\alpha'(q - \alpha p) + \alpha\beta' p' &\geq 0, \\
\alpha(p - \beta q) + \beta\alpha' q' &\geq 0;
\end{align*}
from which you will see that your conditions are narrower than those which the data are really subject to. If your conditions are satisfied, the data will be consistent; but the converse of this proposition does not hold.

\newpage
``4thly.\ You remark that my solution of the problem, in which the independence of $A$ and $B$ is not assumed, but in which the probabilities are otherwise the same as in yours, is only applicable when
\[
\alpha' + \alpha p \geq \beta q, \quad \beta' + \beta q \geq \alpha p;
\]
but you do not appear to have noticed that these are actually the conditions of consistency in the data. Unless these are satisfied, the data cannot possibly be furnished by experience.

``5thly.\ You remark that I have solved the problem under what you call the `concomitance' statement, and not the `causation' statement. I think that every problem stated in the `causation' form admits, if capable of scientific treatment, of reduction to the `concomitance' form. I admit it would have been better, in stating my problem, not to have employed the word `cause' at all. But the introduction of the hypothesis of the independence of $A$ and $B$ does not affect the nature of the problem.

``6thly.\ The $x$, $s$, \&c., about the interpretation of which you inquire, are the probabilities of ideal events in an ideal problem connected by a formal relation with the real one. I should fully concede that the auxiliary probabilities which are employed in my method always refer to an ideal problem; but it is one, the form of which, as given by the calculus of logic, is not arbitrary. Nor does its connexion with the real problem appear to me arbitrary. It involves an extension, but as it seems to me, a perfectly scientific extension, of the principles of the ordinary theory of probabilities. On this subject, however, I have but little to add to what I have said. Transactions of the Royal Society of Edinburgh, vol.\ xxi.\ part 4, ``On the Application of the Theory of Probabilities \&c.''

``7thly.\ The problem, as stated by me, and then modified by the simple introduction of the hypothesis of the independence of $A$ and $B$, must admit of solution by my method; and that solution ought to impose no restriction beyond the conditions of possible experience noted in (M).

I should be extremely glad if mathematicians would examine the analytical questions connected with the application of my method. There can, I think, after the partial proofs which I have given, exist no doubt that the conditions of applicability of the solutions are always identical with the conditions of consistency in the data, i.e.\ with what I have called, in the paper above referred to, the conditions of possible experience. The proof of the general proposition would involve the showing that a certain functional determinant consists solely of positive terms, with some connected theorems which appear to me to be of considerable analytical interest.

``8thly.\ I certainly think your paper deserving of publication. If you think proper to add any or the whole of my remarks, you can do so, with of course any comments of your own.''

\medskip

I remark upon Prof.\ Boole's observations:

\medskip

1st.\ I do assume that the causes $A$ and $B$ are absolutely independent of, and uninfluenced by each other; viz.\ not only the probability of $A$ acting, but also the probability of its acting efficiently, are each of them the same whether $B$ does not act, or acts inefficiently, or acts efficiently; and the like for $B$.

2ndly.\ I do assume that the same experience which establishes
\[
\text{Prob.\ } A = \alpha, \quad \text{Prob.\ } B = \beta,
\]
would in the long run establish
\[
\text{Prob.\ } AB = \alpha\beta;
\]
if it does not, \textit{cadit qu\ae stio}, the causes are not independent.

3rdly.\ I assume not only
\[
\text{Prob.\ } A = \alpha, \quad \text{Prob.\ } B = \beta, \quad \text{Prob.\ } AB = \alpha\beta,
\]
but also as 1st above stated; and I consider that, inasmuch as the result of the investigation is to show that the conditions $q - \alpha p \geq 0$, $p - \beta q \leq 0$ are necessary and sufficient conditions, it is also a result of the investigation that these are the conditions of consistency among the data, viz.\ the conditions in order that the data may be consistent with the above assumptions as to the independence of the causes. It is clear that since, as just stated, I do assume something beyond the last-mentioned three equations, the conditions of consistency ought to be narrower than those in Prof.\ Boole's 3rdly.

4thly.\ I had not overlooked, but I ought to have stated, that Prof.\ Boole's conditions were actually the conditions of consistency in the data.

5thly.\ I contend that the conception of $A$ and $B$ as causes does alter the nature of the problem. For when $A$ and $B$ are conceived of as causes, there is a definite notion of the efficient or inefficient action of $A$ or $B$; and in particular when they both act, one of them, say $A$, may act inefficiently. But according to the concomitance statement, then either there is no such notion as that of the efficient or non-efficient happening of $A$ or $B$ (I believe this to be so), or else the only notion of efficient or inefficient happening is happening in concomitance or in non-concomitance with $E$; but in this view, if $A$, $B$, $E$ all happen, then $A$ and $B$ each of them happens efficiently. The argument is to me conclusive as to the diversity of the two problems.

6thly.\ I do not in anywise assert, or even suppose, that the ideal problem is arbitrary, or that its connexion with the real problem is arbitrary. I simply do not know what the ideal problem is; I do not know the point of view, or the assumed mental state of knowledge or ignorance according to which $x$, $y$, $s$, $t$ are the probabilities of $A$, $B$, $AE$, $BE$. It is to be borne in mind that $x$, $y$, $s$, $t$ are, in Prof.\ Boole's solution, determined as numerical quantities included between the limits 0 and 1, i.e.\ as quantities which are or may be actual probabilities. What I desiderate is, that Prof.\ Boole should give for his auxiliary quantities $x$, $y$, $s$, $t$ such an explanation of the meaning as I have given for my auxiliary quantities $\lambda$, $\mu$. I do not find any such explanation in the memoir referred to.

7thly\ and 8thly.\ No remark is necessary.

\medskip
\noindent March 29, 1862.

\medskip

Prof.\ Boole, in his reply, dated April 2, writes, ``No such explanation as you desiderate of the interpretation of the auxiliary quantities in my method of solution is possible; because they are not of the nature of additional data, and their introduction does not limit the problem as any hypotheses which are of that nature do. I do not see any difficulty whatever in the conception of the ideal problem.''

We thus join issue as follows: Prof.\ Boole says that there is no difficulty in understanding, I say that I do not understand, the rationale of his solution.

It may be remarked that the question may be, not to find any actual probability whatever, but only to find a Boolian probability or probabilities. Thus the equations (L), p.\ 356, omitting the last member, which alone involves $u$, determine in terms of the data $\alpha$, $\beta$, $\alpha p$, $\beta q$ the Boolian probabilities $x$, $y$, $s$, $t$ of the events $A$, $B$, $AE$, $BE$.

In a subsequent hastily-written letter, Prof.\ Boole gives an explanation of the equations (L), which appears to me little more than a translation of these equations into ordinary language.

\medskip
\noindent April 16, 1862.


\newpage
% =====================================================================
% Paper 319 — p. 85
% =====================================================================
\section*{319. Postscript to the Paper ``On a Question in the Theory of Probabilities.''}
\addcontentsline{toc}{section}{319. Postscript to the Paper ``On a Question in the Theory of Probabilities''}

[From the \textit{Philosophical Magazine}, vol.\ xxiii.\ (1862), pp.\ 470--471.]

\medskip

See 318. The present paper reproduces Wilbraham's discussion of Boole's Solution; and concludes with the remark ``Professor Boole wishes me to mention that he has succeeded in obtaining a demonstration of the analytical theorem arising from his theory referred to in his Reply in my paper.'' See ante p.\ 82, 7thly.


\newpage
% =====================================================================
% Paper 320 — pp. 86–88
% =====================================================================
\section*{320. On the Transcendent $\mathrm{gd}\, u = \tfrac{1}{i}\log\tan(\tfrac{1}{4}\pi + \tfrac{1}{2}iu)$.}
\addcontentsline{toc}{section}{320. On the Transcendent gd u = (1/i) log tan(pi/4 + iu/2)}

[From the \textit{Philosophical Magazine}, vol.\ xxiv.\ (1862), pp.\ 19--21.]

\medskip

The elliptic functions which correspond to the modulus $k=1$ reduce themselves, as is well known, to circular functions. The case in question plays implicitly an important part in the general theory, and it has been particularly studied by Gudermann, and by Dr Booth in connexion with his theory of parabolic logarithms. But in the absence of a notation corresponding to that used for elliptic functions in general, the theory has not, it appears to me, been exhibited in its proper form. The defect is very easily supplied: using for $\mathrm{am}\, u$, to the modulus 1, the notation $\mathrm{gd}\, u$ (Gudermannian of $u$), then if
\[
u = \int_0^{\phi}\frac{d\phi}{\cos\phi} = \log\tan\!\left(\tfrac{1}{4}\pi + \tfrac{1}{2}\phi\right),
\]
we have
\[
\phi = \mathrm{gd}\, u;
\]
and hence, observing that the equation between $u$ and $\phi$ is
\[
u = \log\frac{1 + \tan\tfrac{1}{2}\phi}{1 - \tan\tfrac{1}{2}\phi},
\]
or, what is the same thing,
\[
\tan\tfrac{1}{2}\phi = \frac{e^u - 1}{e^u + 1};
\]
and that we have
\begin{align*}
\log\tan\!\left(\tfrac{1}{4}\pi + \tfrac{1}{2}iu\right) &= \log\frac{1 + \tan\tfrac{1}{2}iu}{1 - \tan\tfrac{1}{2}iu} \\
&= \log\frac{\cos\tfrac{1}{2}iu + \sin\tfrac{1}{2}iu}{\cos\tfrac{1}{2}iu - \sin\tfrac{1}{2}iu} \\
&= \log\frac{e^{iu} + e^{-iu} + i(e^{iu} - e^{-iu})}{e^{iu} + e^{-iu} - i(e^{iu} - e^{-iu})} \\
&= \log\frac{e^{iu} + 1 + i(e^{iu} - 1)}{e^{iu} + 1 - i(e^{iu} - 1)} \\
&= \log\frac{1 + i\tan\tfrac{1}{2}\phi}{1 - i\tan\tfrac{1}{2}\phi} = \log e^{i\phi} = i\phi,
\end{align*}
or if
\[
u = \log\tan\!\left(\tfrac{1}{4}\pi + \tfrac{1}{2}\phi\right),
\]
then
\[
\phi = \tfrac{1}{i}\log\tan\!\left(\tfrac{1}{4}\pi + \tfrac{1}{2}ui\right);
\]
and substituting for $\phi$ its value, we obtain
\[
\mathrm{gd}\, u = \tfrac{1}{i}\log\tan\!\left(\tfrac{1}{4}\pi + \tfrac{1}{2}ui\right),
\]
which is the definition of the transcendent $\mathrm{gd}\, u$. It is to be noticed that $\mathrm{gd}\, u$, although exhibited in an imaginary form, is a real function of $u$; and, moreover, that it is an odd function, viz.\ we have
\[
\mathrm{gd}(-u) = -\mathrm{gd}(u),
\]
and therefore also
\[
\mathrm{gd}(0) = 0.
\]

The original equation,
\[
u = \log\tan\!\left(\tfrac{1}{4}\pi + \tfrac{1}{2}\phi\right),
\]
written under the form
\[
u = \tfrac{1}{i}\log\tan\!\left(\tfrac{1}{4}\pi + \tfrac{1}{2}\!\left(\tfrac{\phi}{i}\right)i\right),
\]
shows that we have
\[
u = \tfrac{1}{i}\,\mathrm{gd}\!\left(\tfrac{\phi}{i}\right) = \tfrac{1}{i}\,\mathrm{gd}\!\left(-i\phi\right);
\]
or substituting for $\phi$ its value $\mathrm{gd}\,u$, we have
\[
u = \tfrac{1}{i}\,\mathrm{gd}\!\left(-i\,\mathrm{gd}\,u\right),
\]
which may also be written
\[
iu = \mathrm{gd}\!\left(i\,\mathrm{gd}\,u\right);
\]
so that $\mathrm{gd}\,u$ is a quasi-periodic function of the second order---a property which has not, at least obviously, any analogue in the general theory. We have
\begin{align*}
\cos\,\mathrm{gd}\,u &= \tfrac{1}{2}(e^{i\,\mathrm{gd}\,u} + e^{-i\,\mathrm{gd}\,u}) \\
&= \tfrac{1}{2}\!\left(\sqrt{\frac{1 + \tan\tfrac{1}{2}ui}{1 - \tan\tfrac{1}{2}ui}} + \sqrt{\frac{1 - \tan\tfrac{1}{2}ui}{1 + \tan\tfrac{1}{2}ui}}\right) \\
&= \frac{1 + \tan^2\tfrac{1}{2}ui}{1 - \tan^2\tfrac{1}{2}ui} = \frac{1}{\cos ui};
\end{align*}
and in like manner
\begin{align*}
\sin\,\mathrm{gd}\,u &= \tfrac{1}{2i}(e^{i\,\mathrm{gd}\,u} - e^{-i\,\mathrm{gd}\,u}) \\
&= \tfrac{1}{2i}\!\left(\sqrt{\frac{1 + \tan\tfrac{1}{2}ui}{1 - \tan\tfrac{1}{2}ui}} - \sqrt{\frac{1 - \tan\tfrac{1}{2}ui}{1 + \tan\tfrac{1}{2}ui}}\right) \\
&= \frac{2\tan\tfrac{1}{2}ui}{i(1 - \tan^2\tfrac{1}{2}ui)} = \frac{\sin ui}{i\cos ui};
\end{align*}
or, as these equations may also be written,
\[
\sec\,\mathrm{gd}\,u = \cos ui = \tfrac{1}{2}(e^u + e^{-u}),
\]
\[
\tan\,\mathrm{gd}\,u = \tfrac{1}{i}\sin ui = \tfrac{1}{2}(e^u - e^{-u});
\]
and from these equations we have
\begin{align*}
\sec\,\mathrm{gd}(u+v) &= \sec\,\mathrm{gd}\,u \cdot \sec\,\mathrm{gd}\,v + \tan\,\mathrm{gd}\,u \cdot \tan\,\mathrm{gd}\,v, \\
\tan\,\mathrm{gd}(u+v) &= \tan\,\mathrm{gd}\,u \cdot \sec\,\mathrm{gd}\,v + \tan\,\mathrm{gd}\,v \cdot \sec\,\mathrm{gd}\,u;
\end{align*}
or, what is the same thing,
\begin{align*}
\sin\,\mathrm{gd}(u+v) &= i \cdot \frac{\sin\,\mathrm{gd}\,u + \sin\,\mathrm{gd}\,v}{1 + \sin\,\mathrm{gd}\,u \cdot \sin\,\mathrm{gd}\,v}, \\
\cos\,\mathrm{gd}(u+v) &= \frac{\cos\,\mathrm{gd}\,u \cdot \cos\,\mathrm{gd}\,v}{1 + \sin\,\mathrm{gd}\,u \cdot \sin\,\mathrm{gd}\,v};
\end{align*}
which forms are at once obtainable from the formulae
\begin{align*}
\sin\,\mathrm{am}(u+v) &= \frac{\sin\,\mathrm{am}\,u\,\cos\,\mathrm{am}\,v\,\Delta\,\mathrm{am}\,v + \sin\,\mathrm{am}\,v\,\cos\,\mathrm{am}\,u\,\Delta\,\mathrm{am}\,u}{1 - k^2\sin^2\!\mathrm{am}\,u\,\sin^2\!\mathrm{am}\,v}, \\
\cos\,\mathrm{am}(u+v) &= \frac{\cos\,\mathrm{am}\,u\,\cos\,\mathrm{am}\,v - \sin\,\mathrm{am}\,u\,\sin\,\mathrm{am}\,v\,\Delta\,\mathrm{am}\,u\,\Delta\,\mathrm{am}\,v}{1 - k^2\sin^2\!\mathrm{am}\,u\,\sin^2\!\mathrm{am}\,v}, \\
\Delta\,\mathrm{am}(u+v) &= \frac{\Delta\,\mathrm{am}\,u\,\Delta\,\mathrm{am}\,v - k^2\sin\,\mathrm{am}\,u\,\sin\,\mathrm{am}\,v\,\cos\,\mathrm{am}\,u\,\cos\,\mathrm{am}\,v}{1 - k^2\sin^2\!\mathrm{am}\,u\,\sin^2\!\mathrm{am}\,v},
\end{align*}
observing that for $k=1$ we have $\Delta\,\mathrm{am} = \cos\,\mathrm{am}$, and that the numerators and denominator each of them divide by
\[
1 - \sin\,\mathrm{am}\,u\,\sin\,\mathrm{am}\,v.
\]

\medskip
\noindent 2, Stone Buildings, W.C., May 7, 1862.


\newpage
% =====================================================================
% Paper 321 — p. 89
% =====================================================================
\section*{321. Final Remarks on Mr Jerrard's Theory of Equations of the Fifth Order.}
\addcontentsline{toc}{section}{321. Final Remarks on Mr Jerrard's Theory of Equations of the Fifth Order}

[From the \textit{Philosophical Magazine}, vol.\ xxiv.\ (1862), p.\ 290.]

\medskip

\noindent C.\ V.


\newpage
% =====================================================================
% Paper 322 — pp. 90–94
% =====================================================================
\section*{322. On the Skew Surface of the Third Order.}
\addcontentsline{toc}{section}{322. On the Skew Surface of the Third Order}

[From the \textit{Philosophical Magazine}, vol.\ xxiv.\ (1862), pp.\ 514--519.]

\medskip

The skew surface of the third order, or ``cubic scroll'' (disregarding a certain special form), may be considered as generated by a line which always passes through three directrices; viz., a plane cubic having a node, and two lines, one of them meeting the cubic in the node, the other of them meeting the cubic in an ordinary point. The analytical investigation possesses some interest as an illustration of the analytical theory of skew surfaces in general.

Take for the equation of the cubic
\[
(x^2 + y^2)\alpha\beta - (\alpha^2 + \beta^2)xy = 0,
\]
which belongs to a cubic having a node at the origin, and passing through the point $(\alpha, \beta)$; and for the equations of the two lines
\[
(x - mz = 0, \quad y - nz = 0), \qquad (x - \alpha = 0, \quad y - \beta = 0).
\]
Then, $(X, Y, Z)$ being current coordinates, the equations of the generating line will be
\[
X = x + AZ, \quad Y = y + BZ;
\]
when this meets the line $(X - mZ = 0, Y - nZ = 0)$, we have
\[
mZ = x + AZ, \quad nZ = y + BZ,
\]
and thence
\[
x(n - B) - y(m - A) = 0;
\]
or, what is the same thing,
\[
nx - my - Bx + Ay = 0:
\]
and when it meets the line $(X - \alpha = 0, Y - \beta = 0)$, we have
\[
\alpha = x + AZ, \quad \beta = y + BZ,
\]
and thence
\[
B(x - \alpha) - A(y - \beta) = 0.
\]
We have thus the system of equations
\begin{align*}
(x^2 + y^2)\alpha\beta - (\alpha^2 + \beta^2)xy &= 0, \\
X &= x + AZ, \\
Y &= y + BZ, \\
nx - my - Bx + Ay &= 0, \\
B(x - \alpha) - A(y - \beta) &= 0;
\end{align*}
from which, eliminating $(A, B, x, y)$, we obtain the equation of the surface.

Writing in the last equation
\[
B = s(x - \alpha), \quad A = s(y - \beta)
\]
(values which give $Bx - Ay = -s(\beta x - \alpha y)$), we find
\begin{align*}
X + \alpha s Z &= (1 + sZ)x, \\
Y + \beta s Z &= (1 + sZ)y, \\
(n + \beta s)x - (m + \alpha s)y &= 0;
\end{align*}
whence also
\[
(n + \beta s)(X + \alpha s Z) - (m + \alpha s)(Y + \beta s Z) = 0,
\]
that is
\[
nX - mY + (n\alpha - m\beta)sZ + s(\beta X - \alpha Y) = 0;
\]
or eliminating $s$ from this equation and the two equations
\begin{align*}
x - X + Z(x - \alpha)s &= 0, \\
y - Y + Z(y - \beta)s &= 0,
\end{align*}
we have
\begin{align*}
\{(n\alpha - m\beta)Z + \beta X - \alpha Y\}(x - X) - Z(x - \alpha)(nX - mY) &= 0, \\
\{(n\alpha - m\beta)Z + \beta X - \alpha Y\}(y - Y) - Z(y - \beta)(nX - mY) &= 0;
\end{align*}
these give
\begin{align*}
\Omega x &= X\{(n\alpha - m\beta)Z + \beta X - \alpha Y\} - \alpha Z(nX - mY) \\
&= -mZ\beta X + X(\beta X - \alpha Y) + mZ\alpha Y \\
&= (X - mZ)(\beta X - \alpha Y),
\end{align*}
and
\begin{align*}
\Omega y &= Y\{(n\alpha - m\beta)Z + \beta X - \alpha Y\} - \beta Z(nX - mY) \\
&= nZ\alpha Y + Y(\beta X - \alpha Y) - nZ\beta X \\
&= (Y - nZ)(\beta X - \alpha Y),
\end{align*}
where
\begin{align*}
\Omega &= (n\alpha - m\beta)Z + (\beta X - \alpha Y) - Z(nX - mY) \\
&= \beta(X - mZ) - \alpha(Y - nZ) - Z\{n(X - mZ) - m(Y - nZ)\} \\
&= (\beta - nZ)(X - mZ) - (\alpha - mZ)(Y - nZ);
\end{align*}
so that
\[
x = \frac{(X - mZ)(\beta X - \alpha Y)}{(\beta - nZ)(X - mZ) - (\alpha - mZ)(Y - nZ)},
\]
\[
y = \frac{(Y - nZ)(\beta X - \alpha Y)}{(\beta - nZ)(X - mZ) - (\alpha - mZ)(Y - nZ)},
\]
which equations give the coordinates $(x, y)$ of the point in which the generating line through the point $(X, Y, Z)$ of the surface meets the cubic
\[
(x^2 + y^2)\alpha\beta - (\alpha^2 + \beta^2)xy = 0.
\]
Substituting these values of $(x, y)$ in the equation of the cubic, we obtain the equation
\begin{align*}
&(\alpha^2 + \beta^2)(X - mZ)(Y - nZ)\{(\beta - nZ)(X - mZ) - (\alpha - mZ)(Y - nZ)\} \\
&\quad - \alpha\beta(\beta X - \alpha Y)\{(X - mZ)^2 + (Y - nZ)^2\} = 0;
\end{align*}
or, as it may be written,
\begin{align*}
&(\alpha^2 + \beta^2)(X - mZ)(Y - nZ)\{\beta(X - mZ) - \alpha(Y - nZ)\} \\
&\quad + (\alpha^2 + \beta^2)(X - mZ)(Y - nZ)Z(mY - nX) \\
&\quad - \alpha\beta(\beta X - \alpha Y)\{(X - mZ)^2 + (Y - nZ)^2\} = 0.
\end{align*}
This equation contains, however, the extraneous factor
\[
\beta(X - mZ) - \alpha(Y - nZ),
\]
which, equated to zero, gives the equation of the plane through the node and the line $(x - mz = 0, y - nz = 0)$. In fact, assuming
\begin{align*}
&(\alpha^2 + \beta^2)(X - mZ)(Y - nZ)Z(mY - nX) - \alpha\beta(\beta X - \alpha Y)\{(X - mZ)^2 + (Y - nZ)^2\} \\
&\qquad = \{\beta(X - mZ) - \alpha(Y - nZ)\}\,\Phi(X, Y, Z),
\end{align*}
it will presently be shown that $\Phi$ is an integral function. Hence, omitting the factor in question, we have
\[
(\alpha^2 + \beta^2)(X - mZ)(Y - nZ) + \Phi(X, Y, Z) = 0,
\]
which is the equation of the surface. It only remains to find $\Phi$: writing for this purpose $X + mZ$, $Y + nZ$ in the place of $X$, $Y$ respectively, and putting for a moment
\[
\Phi(X + mZ, Y + nZ, Z) = \Phi',
\]
we have
\[
(\alpha^2 + \beta^2)XYZ(mY - nX) - \alpha\beta\{\beta(X + mZ) - \alpha(Y + nZ)\}(X^2 + Y^2) = (\beta X - \alpha Y)\Phi';
\]
that is
\[
(\beta X - \alpha Y)\Phi' = Z\{(\alpha^2 + \beta^2)XY(mY - nX) - (X^2 + Y^2)\alpha\beta(m\beta - n\alpha)\} - (\beta X - \alpha Y)\alpha\beta(X^2 + Y^2);
\]
or, effecting the division,
\[
\Phi' = Z\{(X^2\alpha - Y^2\beta)(\alpha n - \beta m) - XY(\alpha^2 m + \beta^2 n)\} - \alpha\beta(X^2 + Y^2),
\]
and then writing $X - mZ$, $Y - nZ$ in the place of $X$, $Y$ respectively, we have
\begin{align*}
\Phi(X, Y, Z) &= Z\{((X - mZ)^2\alpha - (Y - nZ)^2\beta)(\alpha n - \beta m) \\
&\quad - (X - mZ)(Y - nZ)(\alpha^2 m + \beta^2 n)\} - \alpha\beta\{(X - mZ)^2 + (Y - nZ)^2\}.
\end{align*}
Hence, finally, the equation of the surface is
\begin{align*}
&(\alpha^2 + \beta^2)(X - mZ)(Y - nZ) - \alpha\beta\{(X - mZ)^2 + (Y - nZ)^2\} \\
&\quad + Z\{((X - mZ)^2\alpha - (Y - nZ)^2\beta)(\alpha n - \beta m) \\
&\quad - (X - mZ)(Y - nZ)(\alpha^2 m + \beta^2 n)\} = 0,
\end{align*}
which is, as it should be, of the third order.

Arranging in powers of $Z$ and reducing, the equation is found to be
\begin{align*}
&(\alpha^2 + \beta^2)XY - \alpha\beta(X^2 + Y^2) \\
&\quad + Z\{-(\alpha^2 + \beta^2)(mY + nX) + (X^2\alpha + Y^2\beta)(m\beta + n\alpha) + \alpha\beta(mX^2 + nY^2) - (\alpha^2 m + \beta^2 n)XY\} \\
&\quad + Z^2\{mn(\alpha^2 + \beta^2 - \alpha^2 X - \beta^2 Y) + (\beta n^2 - \alpha m^2)(\beta X - \alpha Y)\} = 0.
\end{align*}
The first form puts in evidence the nodal line
\[
(X - mZ = 0, \quad Y - nZ = 0),
\]
and the second form puts in evidence the simple line
\[
(X - \alpha = 0, \quad Y - \beta = 0).
\]
But to obtain a more convenient form, write for a moment $X - mZ = P$, $Y - nZ = Q$; the equation is
\[
(\alpha^2 + \beta^2)PQ - \alpha\beta(P^2 + Q^2) + Z\{(P^2\alpha - Q^2\beta)(n\alpha - m\beta) - PQ(m\alpha^2 + n\beta^2)\} = 0,
\]
or, as this may be written,
\[
(\alpha^2 + \beta^2)PQ + (\alpha^2 P - \beta^2 Q)Z(Pn - Qm) + \alpha\beta\{-P^2 - Q^2 - Z(mP^2 + nQ^2)\} = 0;
\]
or, observing that $X = P + mZ$, $Y = Q + nZ$, and thence
\[
PY - QX = Z(Pn - Qm), \quad P^2 X + Q^2 Y = P^3 + Q^3 + Z(mP^2 + nQ^2),
\]
the equation becomes
\[
(\alpha^2 + \beta^2)PQ + (\alpha^2 P - \beta^2 Q)(PY - QX) - \alpha\beta(P^2 X + Q^2 Y) = 0,
\]
or, what is the same thing,
\[
(\alpha P^2 - \beta Q^2)(\alpha Y - \beta X) + PQ(\alpha^2 + \beta^2 - \alpha^2 X - \beta^2 Y) = 0;
\]
whence, making a slight change in the form, and restoring for $P$, $Q$ their values, the equation is
\begin{align*}
&\{\alpha(X - mZ)^2 - \beta(Y - nZ)^2\}\{\alpha(Y - \beta) - \beta(X - \alpha)\} \\
&\quad - (X - mZ)(Y - nZ)\{\alpha^2(X - \alpha) + \beta^2(Y - \beta)\} = 0,
\end{align*}
a form which puts in evidence as well the simple line $(X - \alpha = 0, Y - \beta = 0)$ as the nodal line $(X - mZ = 0, Y - nZ = 0)$.

If $Z = 0$, we have
\[
(\alpha X^2 - \beta Y^2)(\alpha Y - \beta X) - XY\{\alpha^2(X - \alpha) + \beta^2(Y - \beta)\} = 0,
\]
which is in fact the cubic curve $(\alpha^2 + \beta^2)XY - \alpha\beta(X^2 + Y^2) = 0$.

Reverting to a former system of equations
\begin{align*}
nx - my - Bx + Ay &= 0, \\
B(x - \alpha) - A(y - \beta) &= 0,
\end{align*}
or, as these may be written,
\begin{align*}
Bx - Ay &= nx - my, \\
B\alpha - A\beta &= nx - my,
\end{align*}
we find
\begin{align*}
B(\beta x - \alpha y) &= (\beta - y)(nx - my), \\
A(\beta x - \alpha y) &= (\alpha - x)(nx - my);
\end{align*}
so that we have
\[
Y - y = \frac{(\beta - y)(nx - my)}{\beta x - \alpha y}\,Z, \quad X - x = \frac{(\alpha - x)(nx - my)}{\beta x - \alpha y}\,Z,
\]
as the equations of the generating line which passes through the point $(x, y)$ of the cubic curve.

\medskip
\noindent 2, Stone Buildings, W.C., October 28, 1862.


\newpage
% =====================================================================
% Paper 323 — pp. 95–97
% =====================================================================
\section*{323. On a Tactical Theorem Relating to the Triads of Fifteen Things.}
\addcontentsline{toc}{section}{323. On a Tactical Theorem Relating to the Triads of Fifteen Things}

[From the \textit{Philosophical Magazine}, vol.\ xxv.\ (1863), pp.\ 59--61.]

\medskip

The school-girl problem may be stated as follows:---``With 15 things to form 35 triads, involving all the 105 duads, and such that they can be divided into 7 systems, each of 5 triads containing all the 15 things.'' A more simple problem is, ``With 15 things, to form 35 triads involving all the 105 duads.''

In the solution which I formerly gave of the school-girl problem (Phil.\ Mag.\ vol.\ xxxvii.\ 1850, [82]), and which may be presented in the form
\begin{center}
\begin{tabular}{r|ccccccc}
        & $a$ & $b$ & $c$ & $d$ & $e$ & $f$ & $g$ \\ \hline
$abc$   &     &     &     & 35  & 17  & 82  & 64  \\
$ade$   &     & 62  & 84  &     &     & 15  & 37  \\
$afg$   &     & 13  & 47  & 86  & 42  &     & 25  \\
$bdf$   & 47  &     & 16  &     & 38  &     & 25  \\
$beg$   & 58  &     & 23  & 14  &     & 56  & 67  \\
$cdg$   & 12  & 78  &     &     & 56  & 34  &     \\
$cef$   & 36  & 45  &     & 27  &     &     & 18  \\
\end{tabular}
\end{center}
(viz.\ the things being $a$, $b$, $c$, $d$, $e$, $f$, $g$, $1$, $2$, $3$, $4$, $5$, $6$, $7$, $8$, the first pentad of triads is $abc$, $d35$, $e17$, $f82$, $g64$, and so for all the seven pentads of triads), there is obviously a division of the 15 things into $(7 + 8)$ things, viz.\ the 35 triads are composed 7 of them each of 3 out of the 7 things, and the remaining 28 each of 1 out of the 7 things, and 2 out of the 8 things: or attending only to the 8 things, there are 28 triads each of them containing a duad of the 8 things, but there is no triad consisting of 3 of the 8 things. More briefly, we may say that in the system there is an 8 without 3, that is, there are 8 things such that no triad of them occurs in the system.

I believe, but am not sure, that in all the solutions which have been given of the school-girl problem there is an 8 without 3.

Now, considering the more simple problem, there are of course solutions which have an 8 without 3 (since every solution of the school-girl problem is a solution of the more simple problem): but it is very easy to show that there is no solution which has a 9 without 3. I wish to show that there is in every solution at least a 6 without 3. This being so, there will be (if they all exist) 3 classes of solutions, viz.\ those which have at most (1) a 6 without 3, (2) a 7 without 3, (3) an 8 without 3. I believe that the first and second classes exist, as well as the third, which is known to do so.

The proposition to be proved is, that given any system of 35 triads involving all the duads of 15 things; there are always 6 things which are a 6 without 3, that is, they are such that no triad of the 6 things is a triad of the system. This will be the case if it is shown that the number of distinct hexads which can be formed each of them containing a triad of the system is less than
\[
\frac{15 \cdot 14 \cdot 13 \cdot 12 \cdot 11 \cdot 10}{1 \cdot 2 \cdot 3 \cdot 4 \cdot 5 \cdot 6} = 5 \cdot 7 \cdot 11 \cdot 13 = 5005,
\]
the entire number of the hexads of 15 things. Now joining to any triad of the system a triad formed out of the remaining 12 things (there are $\frac{12 \cdot 11 \cdot 10}{1 \cdot 2 \cdot 3} = 4 \cdot 5 \cdot 11 = 220$ such triads), we obtain in all $(220 \times 35 =)\, 7700$ hexads, each of them containing a triad of the system. But these 7700 hexads are not all of them distinct. For, first, considering any triad of the system, there are in the system 16 other triads, each of them having no thing in common with the first-mentioned triad. (In fact if e.g.\ $123$ is a triad of the system, then the system, since it contains all the duads, must have besides 6 triads containing 1, 6 triads containing 2, and 6 triads containing 3, and therefore $35 - 1 - 6 - 6 - 6 = 16$ triads not containing 1, 2, or 3.) Hence we have $\frac{16 \times 35}{2} = 280$ hexads, each of them composed of two triads of the system; and since each of these hexads can be derived from either of its two component triads, these 280 hexads present themselves twice over among the 7700 hexads.

Secondly, there are in the system seven triads containing each of them the same one thing, e.g.
\[
123,\ 145,\ 167,\ 189,\ 1.10.11,\ 1.12.13,\ 1.14.15,
\]
and therefore $\frac{7 \cdot 6}{2} = 21$ pairs such as $123$, $145$ containing the thing $1$, and therefore $(15 \times 21 =)\, 315$ pairs such as $\alpha\beta\gamma$, $\alpha\delta\epsilon$.

\newpage
And for any such pair, combining with $\alpha\beta\gamma\delta\epsilon$ any one of the remaining 10 things, we have 10 hexads $\alpha\beta\gamma\delta\epsilon\zeta$, each of them derivable from either of the triads $\alpha\beta\gamma$, $\alpha\delta\epsilon$, that is, we have $(315 \times 10 =)\, 3150$ hexads which present themselves twice over among the 7700 hexads. The hexads not belonging to one or other of the foregoing classes are derived each of them from a single triad only of the system, and they present themselves once among the 7700 hexads. This number is consequently made up as follows, viz.
\begin{center}
\begin{tabular}{r@{\ hexads each\ }l@{\ \ }r@{\ }r}
$280$  & twice & $=$ & $560$  \\
$3150$ & twice & $=$ & $6300$ \\
$840$  & once  & $=$ & $840$  \\ \cline{4-4}
$4270$ &       &     & $7700$ \\
\end{tabular}
\end{center}
or there are in all 4270 distinct hexads; and since this is less than 5005, it follows that there are hexads not containing any triad of the system: there must in fact be $(5005 - 4270 =)\, 735$ such hexads. The theorem in question is thus shown to be true.

\medskip
\noindent 2, Stone Buildings, W.C., November 24, 1862.

\medskip
\noindent C.\ V.


\newpage
% =====================================================================
% Paper 324 — pp. 98–99
% =====================================================================
\section*{324. Note on a Theorem Relating to Surfaces.}
\addcontentsline{toc}{section}{324. Note on a Theorem Relating to Surfaces}

[From the \textit{Philosophical Magazine}, vol.\ xxv.\ (1863), pp.\ 61, 62.]

\medskip

The following apparently self-evident geometrical theorem requires, I think, a proof; viz.\ the theorem is---``If every plane section of a surface of the order $m+n$ break up into two curves of the orders $m$ and $n$ respectively, then the surface breaks up into two surfaces of the orders $m$, $n$ respectively.''

To fix the ideas, suppose $n = 2$. Imagine any line meeting the surface in $m+2$ points, the section includes a conic which meets the line in two of the $m+2$ points, say the points $A$, $A'^{(1)}$. Suppose that the plane revolves round the line $AA'$, the section will always include a conic which passes through these same two points $A$, $A'$; and it is to be shown that the sheet, the locus of this conic, is a surface of the second order. In fact the conic in question, say $APA'$, by its intersection with an arbitrary plane traces out a branch of the intersection of the given surface with the arbitrary plane. And if $ABA'B'$ be the conic in any particular plane through $A$, $A'$, and if the arbitrary plane meet this conic in the points $B$, $B'$, then the branch passes through these points $B$, $B'$. Imagine the plane $ABA'B'$ revolving round $BB'$ until it coincides with the arbitrary plane; the section includes a conic passing through the points $B$, $B'$, and the before-mentioned branch is this conic; that is, the conic $APA'$ by its intersection with an arbitrary plane traces out a conic; or, what is the same thing, the sheet, the locus of the conic $APA'$, is met by an arbitrary plane in a conic, that is, the sheet is a surface of the second order; and the given surface thus includes a surface of the second order, and is therefore made up of two surfaces of the orders $m$ and $2$ respectively. The demonstration seems to me to add at least something to the evidence of the theorem asserted, but I should be glad if a more simple one could be found. Analytically, the theorem is---``If
\[
(x, y, z, \alpha x + \beta y + \gamma z)^{m+n},
\]
where $(\alpha, \beta, \gamma)$ are arbitrary, break up into factors $(x, y, z)^m$, $(x, y, z)^n$, rational in regard to $(x, y, z)$, then $(x, y, z, w)^{m+n}$ breaks up into factors $(x, y, z, w)^m$, $(x, y, z, w)^n$, rational in regard to $(x, y, z, w)$.'' It would at first sight appear that $(\alpha, \beta, \gamma)$ being arbitrary, these quantities can only enter into the factors of $(x, y, z, \alpha x + \beta y + \gamma z)^{m+n}$ through the quantity $\alpha x + \beta y + \gamma z$; that is, that the factors in question, considered as functions of $(x, y, z, \alpha, \beta, \gamma)$, are of the form
\[
(x, y, z, \alpha x + \beta y + \gamma z)^m, \quad (x, y, z, \alpha x + \beta y + \gamma z)^n;
\]
and then replacing the arbitrary quantity $\alpha x + \beta y + \gamma z$ by $w$, the factors of $(x, y, z, w)^{m+n}$ will be $(x, y, z, w)^m$, $(x, y, z, w)^n$. But the objection proves too much; for in a similar way it would follow that if $(x, y, \alpha x + \beta y)^{m+n}$, where $\alpha$, $\beta$ are arbitrary, breaks up into the factors $(x, y)^m$, $(x, y)^n$, rational in regard to $(x, y)$ (and \textit{qua} homogeneous function of two variables it always does so break up), then $(x, y, z)^{m+n}$ would in like manner break up into the factors $(x, y, z)^m$, $(x, y, z)^n$, rational in regard to $(x, y, z)$: and a simple example will show that it is not true that the factors of $(x, y, \alpha x + \beta y)^{m+n}$ only contain $(\alpha, \beta)$ through $\alpha x + \beta y$; in fact, if the function be $= x^2 + y^2 + (\alpha x + \beta y)^2$, then the factor is
\[
\frac{1}{\sqrt{\alpha^2 + \beta^2 + 1}}\{(\alpha^2 + 1)x + (\alpha\beta + \beta\sqrt{\alpha^2 + \beta^2 + 1})y\},
\]
which cannot be exhibited as a function of $\alpha$, $\beta$, $\alpha x + \beta y$.

I am not acquainted with any analytical demonstration; the geometrical one cannot easily be exhibited in an analytical form.

\medskip
\noindent 2, Stone Buildings, W.C., November 26, 1862.

\bigskip
\noindent\rule{0.3\textwidth}{0.4pt}\\
\footnotesize $^{(1)}$ The figure referred to will be at once understood by considering $A$, $A'$ as the poles of an ellipsoid, or say of a sphere, $ABA'B'$ the meridian of projection, $APA'$ any other meridian, $BPB'$ the equator or any other great circle.


\newpage
% =====================================================================
% Paper 325 — p. 100 (beginning only, continues on p. 101)
% =====================================================================
\section*{325. Note on a Theorem Relating to a Triangle, Line, and Conic.}
\addcontentsline{toc}{section}{325. Note on a Theorem Relating to a Triangle, Line, and Conic}

[From the \textit{Philosophical Magazine}, vol.\ xxv.\ (1863), pp.\ 181--183.]

\medskip

I \textsc{find}, among my papers headed ``Generalization of a Theorem of Steiner's,'' an investigation leading to the following theorem, viz.:

Consider a triangle, a line, and a conic; with each vertex of the triangle join the point of intersection of the line with the polar of the same vertex in regard to the conic; in order that the three joining lines may meet in a point, the line must be a tangent to a curve of the third class; if, however, the conic break up into a pair of lines, or in a certain other case, the curve of the third class will break up into a point, and a conic inscribed in the triangle.

Let the equations of the sides of the triangle be
\[
x = 0, \quad y = 0, \quad z = 0,
\]
the equation of the conic
\[
(a, b, c, f, g, h \,\vert\, x, y, z)^2 = 0,
\]
and that of the line
\[
\lambda x + \mu y + \nu z = 0;
\]
then the polar of the vertex $(y = 0, z = 0)$ has for its equation
\[
ax + hy + gz = 0;
\]
it therefore meets the line $\lambda x + \mu y + \nu z = 0$ in the point
\[
x : y : z = h\nu - g\mu : g\lambda - a\nu : a\mu - h\lambda.
\]

\noindent[Continues on p.~101.]


\end{document}
